\subsection{为什么首轮统计量不是捷径}\label{sec:counterfactual}

轮内统计量能在产生它们的译码器内部预测失败~\cite{kim2021logistic,paiva2022pickyfix}。我们检验一个在参考变体首轮计数上训练的代理模型能否迁移到另一套调度。模型是对正确翻转数与错误翻转数做逻辑回归拟合，在 46 把密钥上训练、在 46 把不相交密钥上评估，每把 3,000 个错误。候选特征按验证集 C0 对数损失选择；预测 RR 为预测的 C0 DFR 除以预测的干预 DFR。在提供标签的臂内，它近乎完美地排序实例（AUC 0.995）；该 AUC 在用于模型选择的验证集上测得，不是最终留出测试。跨调度时它失败了：对留出的偏移 $+1$、$+2$、$+3$，它预测 RR 0.79、0.16、0.04，而配对测量为 3.46、3.86、3.97。两个计数是处理后变量：正确翻转数偏低在 C0 下标记困难实例，在抬高阈值下却标记被推迟的翻转。因此这个特定的代理模型是产生它的调度的诊断工具，而不是另一套调度的选择器。我们不主张所有首轮统计量都失败。这正是本文直接配对测量、而不是做预测的原因。
